% !TEX root = ../main.tex
% ---------------------------------------------------------------
% GENERATED FILE - DO NOT EDIT.
% Produced by docs/tools/render_reference.py from the repository source.
% Regenerate with: make docs
% ---------------------------------------------------------------

\apibreak
\section{\texttt{QdrantVectorStore}}
\label{class:feedback_intelligence_agent.qdrant_store.QdrantVectorStore}\index{Classes!QdrantVectorStore}

Cosine-similarity vector store backed by a Qdrant collection.
\apifield{Usage}
\begin{lstlisting}[style=usage, language=Python]
QdrantVectorStore(self, dim: int, *, url: str = 'http://localhost:6333', collection_name: str = 'feedback_intelligence', client: QdrantClient | None = None) -> None
\end{lstlisting}
\apifield{Arguments}
\begin{arglist}
  \item[\texttt{dim}] \texttt{int}, required.
  \item[\texttt{url}] \texttt{str}, default \texttt{'http:/\allowbreak{}/\allowbreak{}localhost:6333'}, keyword-only.
  \item[\texttt{collection\_\allowbreak{}name}] \texttt{str}, default \texttt{'feedback\_\allowbreak{}intelligence'}, keyword-only.
  \item[\texttt{client}] \texttt{QdrantClient \textbar{} None}, default \texttt{None}, keyword-only.
\end{arglist}
\apifield{Raises}\texttt{ValueError}, \texttt{QdrantStoreError}
\apifield{Properties}
\begin{arglist}
  \item[\texttt{size}] Number of indexed chunks in the collection.
  \item[\texttt{chunks}] Return all indexed chunks, reconstructed from point payloads.
\end{arglist}
\apifield{Details}The collection is created on demand using \texttt{Distance.\allowbreak{}COSINE} so the scoring orientation matches the in-memory store: higher scores mean more similar.
\apifield{Source}\texttt{src/\allowbreak{}feedback\_\allowbreak{}intelligence\_\allowbreak{}agent/\allowbreak{}qdrant\_\allowbreak{}store.\allowbreak{}py:56}~(\href{https://github.com/DiogoRibeiro7/feedback-intelligence-agent/blob/b6cbc710ffdfa5f26f700c999d594c67588c1e16/src/feedback_intelligence_agent/qdrant_store.py#L56}{online}), class, module \cref{module:feedback_intelligence_agent.qdrant_store}

\subsection{\texttt{add}}
\label{method:feedback_intelligence_agent.qdrant_store.QdrantVectorStore.add}\index{Methods!add}
Upsert chunks and their vectors into the collection.
\apifield{Usage}
\begin{lstlisting}[style=usage, language=Python]
add(self, chunks: list[DocumentChunk], vectors: npt.NDArray[np.float64]) -> None
\end{lstlisting}
\apifield{Arguments}
\begin{arglist}
  \item[\texttt{chunks}] \texttt{list[DocumentChunk]\allowbreak{}}, required.
  \item[\texttt{vectors}] \texttt{npt.\allowbreak{}NDArray[np.\allowbreak{}float64]\allowbreak{}}, required.
\end{arglist}
\apifield{Raises}\texttt{ValueError}
\apifield{Source}\texttt{src/\allowbreak{}feedback\_\allowbreak{}intelligence\_\allowbreak{}agent/\allowbreak{}qdrant\_\allowbreak{}store.\allowbreak{}py:129}~(\href{https://github.com/DiogoRibeiro7/feedback-intelligence-agent/blob/b6cbc710ffdfa5f26f700c999d594c67588c1e16/src/feedback_intelligence_agent/qdrant_store.py#L129}{online})

\subsection{\texttt{search}}
\label{method:feedback_intelligence_agent.qdrant_store.QdrantVectorStore.search}\index{Methods!search}
Return the most similar chunks for a query vector (highest score first).
\apifield{Usage}
\begin{lstlisting}[style=usage, language=Python]
search(self, query_vector: npt.NDArray[np.float64], *, top_k: int = 4) -> list[SearchResult]
\end{lstlisting}
\apifield{Arguments}
\begin{arglist}
  \item[\texttt{query\_\allowbreak{}vector}] \texttt{npt.\allowbreak{}NDArray[np.\allowbreak{}float64]\allowbreak{}}, required.
  \item[\texttt{top\_\allowbreak{}k}] \texttt{int}, default \texttt{4}, keyword-only.
\end{arglist}
\apifield{Value}\texttt{list[SearchResult]\allowbreak{}}
\apifield{Raises}\texttt{ValueError}
\apifield{Source}\texttt{src/\allowbreak{}feedback\_\allowbreak{}intelligence\_\allowbreak{}agent/\allowbreak{}qdrant\_\allowbreak{}store.\allowbreak{}py:151}~(\href{https://github.com/DiogoRibeiro7/feedback-intelligence-agent/blob/b6cbc710ffdfa5f26f700c999d594c67588c1e16/src/feedback_intelligence_agent/qdrant_store.py#L151}{online})
